\section*{結言}
今回の観測では、AIが実世界の設計作業を支援する具体例と、その推論結果の出所を検証する技術が並んで注目された。設計の反復速度を高めることと、生成過程への信頼を確保することは、AIを実務に組み込むうえで相補的な課題である。

一方、計算基盤への投資拡大と、従業員の発言、金融口座へのアクセスといった統制上の問題が同時に浮上した。AIエージェントの利用範囲が広がるほど、権限・責任・検証の設計が重要になることを示している。

次期モデルの噂や法律特化モデルの性能など、裏付けが十分でない情報も含まれる。Grokの観測窓における話題と、各社が公式に確認した事実とを混同せず、今後の続報を追う必要がある。
